% --- Archon physics formalization source begin ---
% archon:physics
% archon:covers PhyXMiniProblems/problem_phyx_mini_0762.lean
% archon:source-report reports/phyx_mini/problem_phyx_mini_0762.source.json

\chapter{Physics problem phyx\_mini\_0762}
\label{ch:PhyXMiniProblems_problem_phyx_mini_0762}

\paragraph{Problem source.}
\#\# Physical scenario

In figure, elevator cabs \$A\$ and \$B\$ are connected by a short cable and can be pulled upward or lowered by the cable above cab \$A\$. Cab \$A\$ has mass \$1700\textbackslash{} kg\$; cab \$B\$ has mass \$1300\textbackslash{} kg\$. A \$12.0\textbackslash{} kg\$ box of catnip lies on the floor of cab \$A\$. The tension in the cable connecting the cabs is \$1.91 \textbackslash{}times 10\textasciicircum{}4\textbackslash{} N\$.

\#\# Question

What is the magnitude of the normal force on the box from the floor?

\#\# Answer choices

- A: A: 168 N

- B: B: 172 N

- C: C: 176 N

- D: D: 180 N

\#\# Figure caption (auxiliary; use the image as primary evidence)

The image contains two sections labeled "A" and "B." Each section appears to depict a rectangular box or container with a metallic outline and is divided into two compartments. 

- **Section A (top):** 
  - The top compartment is filled with a light green color.
  - Below it, there is a thinner compartment filled with a slightly darker green color.
  - A yellow rectangle is placed in the lower part of the dark green compartment.
  - A cable or wire is attached to the top center of the container.

- **Section B (bottom):**
  - Mirrors the structure of Section A.
  - The top compartment is light green.
  - Below it is a narrower compartment with darker green.
  - There is no additional object within the compartments in this section.
  - A cable or wire is attached to the top center of the container.

The sections seem to represent some kind of mechanical or fluid system with distinct areas, represented by the colored compartments.

\#\# Dataset metadata

Dynamics; Physical Model Grounding Reasoning, Multi-Formula Reasoning

\paragraph{Recorded answer/context.}
C: C: 176 N

\paragraph{Figure/image path.}
/root/proposal\_for\_physic/science-mango/phyx\_mini\_run/phyx\_data/test\_image/762.png

\paragraph{Formalization target.}
create a compiling Lean file with sorry bodies at `PhyXMiniProblems/problem\_phyx\_mini\_0762.lean`. The Lean declarations must preserve the physical quantities, dimensions or dimensional roles, figure labels, governing-law hypotheses, and final relation expressed by this problem.
Use Mathlib/Physlib names found through LeanExplore where available. If a physics API is missing, introduce faithful local abstractions rather than scalar placeholder aliases.

\begin{theorem}[Physics formalization target]
\label{thm:physics:phyx_mini_0762:target}
\lean{PhyXMiniProblems.ProblemPhyXMini0762.problem_phyx_mini_0762}
\uses{def:physics:phyx-mini-0762:phyxminiproblems-problemphyxmini0762-massinkilograms, def:physics:phyx-mini-0762:phyxminiproblems-problemphyxmini0762-forceinnewtons, def:physics:phyx-mini-0762:phyxminiproblems-problemphyxmini0762-connectedelevatorsetup, def:physics:phyx-mini-0762:phyxminiproblems-problemphyxmini0762-matchesconnectedelevatorfigure, def:physics:phyx-mini-0762:phyxminiproblems-problemphyxmini0762-hasconnectedelevatorproblemdata, def:physics:phyx-mini-0762:phyxminiproblems-problemphyxmini0762-satisfiesconnectedelevatordynamics, def:physics:phyx-mini-0762:phyxminiproblems-problemphyxmini0762-normalforceanswerchoice, def:physics:phyx-mini-0762:phyxminiproblems-problemphyxmini0762-displayednormalforceinnewtons, def:physics:phyx-mini-0762:phyxminiproblems-problemphyxmini0762-roundstonearestnewton}
The assigned autoformalize agent should translate this physics problem into Lean declarations in the covered file, with theorem and lemma proof bodies written as `by sorry`.
\end{theorem}
\begin{proof}
This is an autoformalization task, not a proof task. Produce statements that can later be proved by the physics prover without weakening the physical model.
\end{proof}

% --- Archon declaration topology begin ---
\section{Declaration topology}
\begin{definition}[acceleration Dimension]
  \label{def:physics:phyx-mini-0762:phyxminiproblems-problemphyxmini0762-accelerationdimension}
  \lean{PhyXMiniProblems.ProblemPhyXMini0762.accelerationDimension}
  The physical dimension of acceleration, L T⁻²
\end{definition}
\begin{proof}
  This is the declared model definition of acceleration Dimension.
\end{proof}
\begin{definition}[force Dimension]
  \label{def:physics:phyx-mini-0762:phyxminiproblems-problemphyxmini0762-forcedimension}
  \lean{PhyXMiniProblems.ProblemPhyXMini0762.forceDimension}
  \uses{def:physics:phyx-mini-0762:phyxminiproblems-problemphyxmini0762-accelerationdimension}
  The physical dimension of force, M L T⁻²
\end{definition}
\begin{proof}
  This is the declared model definition of force Dimension.
\end{proof}
\begin{definition}[Mass Quantity]
  \label{def:physics:phyx-mini-0762:phyxminiproblems-problemphyxmini0762-massquantity}
  \lean{PhyXMiniProblems.ProblemPhyXMini0762.MassQuantity}
  A nonnegative, unit-independent physical mass
\end{definition}
\begin{proof}
  This is the declared model definition of Mass Quantity.
\end{proof}
\begin{definition}[Acceleration Magnitude]
  \label{def:physics:phyx-mini-0762:phyxminiproblems-problemphyxmini0762-accelerationmagnitude}
  \lean{PhyXMiniProblems.ProblemPhyXMini0762.AccelerationMagnitude}
  \uses{def:physics:phyx-mini-0762:phyxminiproblems-problemphyxmini0762-accelerationdimension}
  A nonnegative, unit-independent acceleration magnitude
\end{definition}
\begin{proof}
  This is the declared model definition of Acceleration Magnitude.
\end{proof}
\begin{definition}[Signed Acceleration]
  \label{def:physics:phyx-mini-0762:phyxminiproblems-problemphyxmini0762-signedacceleration}
  \lean{PhyXMiniProblems.ProblemPhyXMini0762.SignedAcceleration}
  \uses{def:physics:phyx-mini-0762:phyxminiproblems-problemphyxmini0762-accelerationdimension}
  A unit-independent signed acceleration along the vertical axis
\end{definition}
\begin{proof}
  This is the declared model definition of Signed Acceleration.
\end{proof}
\begin{definition}[Force Magnitude]
  \label{def:physics:phyx-mini-0762:phyxminiproblems-problemphyxmini0762-forcemagnitude}
  \lean{PhyXMiniProblems.ProblemPhyXMini0762.ForceMagnitude}
  \uses{def:physics:phyx-mini-0762:phyxminiproblems-problemphyxmini0762-forcedimension}
  A nonnegative, unit-independent force magnitude
\end{definition}
\begin{proof}
  This is the declared model definition of Force Magnitude.
\end{proof}
\begin{definition}[mass In Kilograms]
  \label{def:physics:phyx-mini-0762:phyxminiproblems-problemphyxmini0762-massinkilograms}
  \lean{PhyXMiniProblems.ProblemPhyXMini0762.massInKilograms}
  \uses{def:physics:phyx-mini-0762:phyxminiproblems-problemphyxmini0762-massquantity}
  Coherent-SI kilogram readout of a physical mass
\end{definition}
\begin{proof}
  This is the declared model definition of mass In Kilograms.
\end{proof}
\begin{definition}[acceleration Magnitude In Meters Per Second Squared]
  \label{def:physics:phyx-mini-0762:phyxminiproblems-problemphyxmini0762-accelerationmagnitudeinmeterspersecondsquared}
  \lean{PhyXMiniProblems.ProblemPhyXMini0762.accelerationMagnitudeInMetersPerSecondSquared}
  \uses{def:physics:phyx-mini-0762:phyxminiproblems-problemphyxmini0762-accelerationmagnitude}
  Coherent-SI metre-per-second-squared readout of an acceleration magnitude
\end{definition}
\begin{proof}
  This is the declared model definition of acceleration Magnitude In Meters Per Second Squared.
\end{proof}
\begin{definition}[signed Acceleration In Meters Per Second Squared]
  \label{def:physics:phyx-mini-0762:phyxminiproblems-problemphyxmini0762-signedaccelerationinmeterspersecondsquared}
  \lean{PhyXMiniProblems.ProblemPhyXMini0762.signedAccelerationInMetersPerSecondSquared}
  \uses{def:physics:phyx-mini-0762:phyxminiproblems-problemphyxmini0762-signedacceleration}
  Coherent-SI signed metre-per-second-squared acceleration component
\end{definition}
\begin{proof}
  This is the declared model definition of signed Acceleration In Meters Per Second Squared.
\end{proof}
\begin{definition}[force In Newtons]
  \label{def:physics:phyx-mini-0762:phyxminiproblems-problemphyxmini0762-forceinnewtons}
  \lean{PhyXMiniProblems.ProblemPhyXMini0762.forceInNewtons}
  \uses{def:physics:phyx-mini-0762:phyxminiproblems-problemphyxmini0762-forcemagnitude}
  Coherent-SI newton readout of a force magnitude
\end{definition}
\begin{proof}
  This is the declared model definition of force In Newtons.
\end{proof}
\begin{definition}[weight In Newtons]
  \label{def:physics:phyx-mini-0762:phyxminiproblems-problemphyxmini0762-weightinnewtons}
  \lean{PhyXMiniProblems.ProblemPhyXMini0762.weightInNewtons}
  \uses{def:physics:phyx-mini-0762:phyxminiproblems-problemphyxmini0762-massquantity, def:physics:phyx-mini-0762:phyxminiproblems-problemphyxmini0762-accelerationmagnitude, def:physics:phyx-mini-0762:phyxminiproblems-problemphyxmini0762-massinkilograms, def:physics:phyx-mini-0762:phyxminiproblems-problemphyxmini0762-accelerationmagnitudeinmeterspersecondsquared}
  The gravitational weight magnitude m g, read in newtons
\end{definition}
\begin{proof}
  This is the declared model definition of weight In Newtons.
\end{proof}
\begin{definition}[Cab Label]
  \label{def:physics:phyx-mini-0762:phyxminiproblems-problemphyxmini0762-cablabel}
  \lean{PhyXMiniProblems.ProblemPhyXMini0762.CabLabel}
  The two cab labels printed beside the primary image
\end{definition}
\begin{proof}
  This is the declared model definition of Cab Label.
\end{proof}
\begin{definition}[Vertical Positive Direction]
  \label{def:physics:phyx-mini-0762:phyxminiproblems-problemphyxmini0762-verticalpositivedirection}
  \lean{PhyXMiniProblems.ProblemPhyXMini0762.VerticalPositiveDirection}
  Sign convention for the scalar vertical acceleration components
\end{definition}
\begin{proof}
  This is the declared model definition of Vertical Positive Direction.
\end{proof}
\begin{definition}[Elevator Figure]
  \label{def:physics:phyx-mini-0762:phyxminiproblems-problemphyxmini0762-elevatorfigure}
  \lean{PhyXMiniProblems.ProblemPhyXMini0762.ElevatorFigure}
  \uses{def:physics:phyx-mini-0762:phyxminiproblems-problemphyxmini0762-cablabel}
  Qualitative information visible in the source image and stated in the scenario. This record contains no mass, force, or acceleration readout
\end{definition}
\begin{proof}
  This is the declared model definition of Elevator Figure.
\end{proof}
\begin{definition}[Connected Elevator Setup]
  \label{def:physics:phyx-mini-0762:phyxminiproblems-problemphyxmini0762-connectedelevatorsetup}
  \lean{PhyXMiniProblems.ProblemPhyXMini0762.ConnectedElevatorSetup}
  \uses{def:physics:phyx-mini-0762:phyxminiproblems-problemphyxmini0762-massquantity, def:physics:phyx-mini-0762:phyxminiproblems-problemphyxmini0762-accelerationmagnitude, def:physics:phyx-mini-0762:phyxminiproblems-problemphyxmini0762-signedacceleration, def:physics:phyx-mini-0762:phyxminiproblems-problemphyxmini0762-forcemagnitude, def:physics:phyx-mini-0762:phyxminiproblems-problemphyxmini0762-cablabel, def:physics:phyx-mini-0762:phyxminiproblems-problemphyxmini0762-verticalpositivedirection, def:physics:phyx-mini-0762:phyxminiproblems-problemphyxmini0762-elevatorfigure}
  The independent physical quantities of the elevator system. In particular, the floor normal force is an observable field; it is not defined from an answer choice or from the requested numerical result
\end{definition}
\begin{proof}
  This is the declared model definition of Connected Elevator Setup.
\end{proof}
\begin{definition}[Matches Connected Elevator Figure]
  \label{def:physics:phyx-mini-0762:phyxminiproblems-problemphyxmini0762-matchesconnectedelevatorfigure}
  \lean{PhyXMiniProblems.ProblemPhyXMini0762.MatchesConnectedElevatorFigure}
  \uses{def:physics:phyx-mini-0762:phyxminiproblems-problemphyxmini0762-connectedelevatorsetup}
  Transcription of the primary image: cab A is the upper cab, cab B is the lower cab, the box is on A's floor, and the cables are vertical
\end{definition}
\begin{proof}
  This is the declared model definition of Matches Connected Elevator Figure.
\end{proof}
\begin{definition}[Has Connected Elevator Problem Data]
  \label{def:physics:phyx-mini-0762:phyxminiproblems-problemphyxmini0762-hasconnectedelevatorproblemdata}
  \lean{PhyXMiniProblems.ProblemPhyXMini0762.HasConnectedElevatorProblemData}
  \uses{def:physics:phyx-mini-0762:phyxminiproblems-problemphyxmini0762-massinkilograms, def:physics:phyx-mini-0762:phyxminiproblems-problemphyxmini0762-forceinnewtons, def:physics:phyx-mini-0762:phyxminiproblems-problemphyxmini0762-connectedelevatorsetup}
  The masses and inter-cab cable tension stated in the problem. No normal-force value or answer label occurs in these independent data
\end{definition}
\begin{proof}
  This is the declared model definition of Has Connected Elevator Problem Data.
\end{proof}
\begin{definition}[Satisfies Connected Elevator Dynamics]
  \label{def:physics:phyx-mini-0762:phyxminiproblems-problemphyxmini0762-satisfiesconnectedelevatordynamics}
  \lean{PhyXMiniProblems.ProblemPhyXMini0762.SatisfiesConnectedElevatorDynamics}
  \uses{def:physics:phyx-mini-0762:phyxminiproblems-problemphyxmini0762-massinkilograms, def:physics:phyx-mini-0762:phyxminiproblems-problemphyxmini0762-signedaccelerationinmeterspersecondsquared, def:physics:phyx-mini-0762:phyxminiproblems-problemphyxmini0762-forceinnewtons, def:physics:phyx-mini-0762:phyxminiproblems-problemphyxmini0762-weightinnewtons, def:physics:phyx-mini-0762:phyxminiproblems-problemphyxmini0762-connectedelevatorsetup}
  The kinematic constraints and the two vertical Newton-second-law equations used in the calculation. The taut short cable makes the two cabs share a vertical acceleration, and a box maintaining contact with A's floor shares cab A's acceleration. On cab B, the inter-cab tension acts upward and weight acts downward. On the box, the floor normal force acts upward and weight acts downward. These are general governing relations and contain no requested numerical normal force
\end{definition}
\begin{proof}
  This is the declared model definition of Satisfies Connected Elevator Dynamics.
\end{proof}
\begin{definition}[Normal Force Answer Choice]
  \label{def:physics:phyx-mini-0762:phyxminiproblems-problemphyxmini0762-normalforceanswerchoice}
  \lean{PhyXMiniProblems.ProblemPhyXMini0762.NormalForceAnswerChoice}
  Labels of the four normal-force answer choices
\end{definition}
\begin{proof}
  This is the declared model definition of Normal Force Answer Choice.
\end{proof}
\begin{definition}[displayed Normal Force In Newtons]
  \label{def:physics:phyx-mini-0762:phyxminiproblems-problemphyxmini0762-displayednormalforceinnewtons}
  \lean{PhyXMiniProblems.ProblemPhyXMini0762.displayedNormalForceInNewtons}
  \uses{def:physics:phyx-mini-0762:phyxminiproblems-problemphyxmini0762-normalforceanswerchoice}
  Whole-newton value printed beside each answer label
\end{definition}
\begin{proof}
  This is the declared model definition of displayed Normal Force In Newtons.
\end{proof}
\begin{definition}[Rounds To Nearest Newton]
  \label{def:physics:phyx-mini-0762:phyxminiproblems-problemphyxmini0762-roundstonearestnewton}
  \lean{PhyXMiniProblems.ProblemPhyXMini0762.RoundsToNearestNewton}
  \uses{def:physics:phyx-mini-0762:phyxminiproblems-problemphyxmini0762-forcemagnitude, def:physics:phyx-mini-0762:phyxminiproblems-problemphyxmini0762-forceinnewtons}
  A physical force magnitude rounds to a displayed whole number of newtons
\end{definition}
\begin{proof}
  This is the declared model definition of Rounds To Nearest Newton.
\end{proof}
% --- Archon declaration topology end ---
% --- Archon physics formalization source end ---
